\newcommand\odp[1]{\par\smallskip{\color{pomd}\bfseries #1}\par}
\odp{1. Wielokrotności}
Ćw. A: 12, 18, 24. Ćw. B: a) np. 0, 7, 14, 21, 28; b) 0, 14, 38, 150. Ćw. C: zielone 36 i 60 cm, czerwone 42 i 56 cm; inne – wielokrotności 12 i 14. Ćw. D: NWW(4, 10) = 20.\par
1. np. a) 9, 18, 27 b) 25, 50, 75 c) 120, 240, 360 d) 3000, 6000, 9000.
2. 1, 2, 3, 6, 9, 18.
3. a) 0, 18, 24, 51, 57, 72, 96 b) 0, 24, 72, 96 c) 0, 34, 51, 85.
4. a) 320 b) 280.
5. a) A = 48, C = 72 b) D = 261, F = 279, G = 288.
6. a) 20 b) 9.
7. P, F, P, P, F.
8. 4 (342 = 38 · 9).
9. a) 12 b) 24 c) 20 d) 347 e) 91 f) 30 g) 30 h) 12.\par
Super zagadka: kreski spotykają się co NWW(25, 20) = 100 cm, więc podwórko ma co najmniej 18 m – tak.

\odp{2. Dzielniki}
Ćw. A: a) 1, 2, 3, 4, 6, 8, 12, 24 b) 1, 3, 5, 15. Ćw. B: 1, 2, 4, 5, 10, 20. Ćw. C: a) 1, 2, 4, 7, 14, 28 b) 1; 2. Ćw. D: b) 1, 2, 3, 4, 6, 9, 12, 18, 36. Ćw. E: 4.\par
1. a) 1, 7 b) 1, 2, 4, 8 c) 1, 2, 4, 8, 16 d) 1, 3, 9, 27 e) 1, 2, 4, 5, 8, 10, 20, 40 f) 1, 2, 4, 8, 16, 32, 64.
2. P, P, F, P, F.
3. 1, 15, 30, 75, 150.
4. Drzwi 16 (1 + 2 + 4 + 8 + 16 = 31).
5. a) 200, 400, 600, 800 b) 150, 300, 450, 600, 750, 900 c) 250, 500, 750 d) 300, 600, 900.
6. a) 22 b) 36.
7. a) 1 + 2 + 3 = 6 – tak b) tak (1 + 2 + 4 + 7 + 14 = 28) c) nie (1 + 2 + 3 + 4 + 6 = 16).
8. np. 150; najmniejsza 120.
9. a) 6 b) 1 c) 5 d) 1 e) 1 f) 3 g) 1 h) 13 i) 1 j) 20 k) 1 l) 5.
10. 12 zestawów (po 2 jabłka i 3 gruszki).
11. a) 13, 26 b) 8, 16, 24 c) 5, 15, 25 d) np. 9 i 18, 9 i 27, 18 i 27.\par
Super zagadka: 21 (1 + 3 + 7 + 21 = 32).

\odp{3. Cechy podzielności}
Ćw. A: a) nie b) tak. Ćw. B: a) 340, 72, 64, 250, 900, 3400 b) 340, 55, 35, 105, 250, 175, 900, 815, 3400 c) 340, 250, 900, 3400 d) 900, 3400. Ćw. C: a) tak b) nie. Ćw. D: 436 i 9612. Ćw. E: nikt (np. 12). Ćw. F: 12, 27, 33, 51, 213 (obie grupy takie same). Ćw. G: 36 – tak, 9, tak; 47 – nie, 11, nie; 58 – nie, 13, nie; 702 – tak, 9, tak; 819 – tak, 18, tak; 974 – nie, 20, nie; 4042 – nie, 10, nie.\par
1. a) 20, 64, 488, 730, 1250, 2006, 4700, 98\,300, 2\,000\,500 b) 20, 85, 315, 730, 1250, 4700, 6205, 98\,300, 2\,000\,500 c) 20, 730, 1250, 4700, 98\,300, 2\,000\,500 d) 4700, 98\,300, 2\,000\,500.
2. Dwie ostatnie cyfry: 00, 25, 50 lub 75.
3. 1324, 516, 892, 21\,436.
4. 724, 936.
5. b) nie; 2028 i 2032 c) 1704, 1948, 2000, 2024 d) 1800, 1900, 2000 – przestępny tylko 2000.
6. 29 lutego 2028 r.
7. a) 135, 30\,012, 8127, 41\,022 b) 135, 8127, 41\,022.
8. a) np. 123 b) nie.
9. 7380.
10. przez 3: A, B, C, D, E; przez 9: B, C.
11. 1\,023\,456\,789 i 9\,876\,543\,210 – tak (suma cyfr 45).
12. a) 9 b) 11.
13. a) KOT, LOT.
14. tak; nie.
15. np. a) 1020, 10\,020 b) 1002, 10\,002 c) 1002, 10\,002.
16. a) np. 105 – tak b) 75 i 4605.
17. a) 2460, 1356, 91\,284 b) 2460, 1356, 91\,284 c) 2460, 74\,020.
18. przez 5: 2, 4, 1, 4, 0; przez 4: 1, 0, 2, 3, 2.
19. a) 1, 2, 0.\par
Super zagadka: 5 (2558 = 12 · 213 + 2).

\odp{4. Liczby pierwsze i złożone}
Ćw. A: 4, 2, 4, 2. Sito: 25 liczb.\par
1. np. 2316 – 2, 4785 – 5, 1818 – 2, 2721 – 3, 7413 – 3.
2. a) 101 i 997 b) 168 c) 830.
3. a) 1, 2003 b) 0, 2, 4, 5, 6, 8.
4. P, P, F, F, P.
5. 1.
6. 103, 827, 1993, 2027.\par
Super zagadka: 3; 4.

\odp{5. Rozkład na czynniki}
1. 18 = 2 · 3 · 3; 45 = 3 · 3 · 5; 72 = 2 · 2 · 2 · 3 · 3; 100 = 2 · 2 · 5 · 5; 126 = 2 · 3 · 3 · 7; 210 = 2 · 3 · 5 · 7.
2. a) 90 b) 56 c) 165.
3. 6; 36; 15; 90.
4. 3; 3.

\textbf{\large Sprawdziany}\par\smallskip
\textbf{\color{pomd}Sprawdzian 1: Wielokrotności}\par
1. 6 {\color{szary}– 962 = 26 · 37.}\par
2. 30 {\color{szary}– 30 to najmniejsza liczba (różna od 0), która dzieli się przez 6, 5, 10.}\par
3. 30 {\color{szary}– Wypisujemy wielokrotności większej liczby: 15, 30 – pierwsza, która dzieli się przez 6, to \textbf{30}.}\par
4. 63 {\color{szary}– 9 · 14 = 126, a 126 : 2 = 63.}\par
5. 81 {\color{szary}– Kolejne wielokrotności liczby 9 różnią się o 9: A = 81, B = 90, C = 99.}\par
6. A. Prawda {\color{szary}– 4 · n = 2 · (2 · n).}\par
7. D. 156 {\color{szary}– 156 = 12 · 13. Pozostałe liczby nie dzielą się przez 12 bez reszty.}\par
8. 135 {\color{szary}– 135 = 45 · 3; następna wielokrotność 180 jest już $\geq$ 162.}\par
9. 2 {\color{szary}– 4 · x − 2 · x = 2 · x = 4, więc x = 4 : 2 = 2.}\par
10. 150 {\color{szary}– 150 = 50 · 3, a poprzednia wielokrotność 100 nie jest większa od 135.}\par
\medskip
\textbf{\color{pomd}Sprawdzian 2: Dzielniki}\par
1. 6 {\color{szary}– Dzielniki 28: 1, 2, 4, 7, 14, 28 – razem 6.}\par
2. 120 {\color{szary}– Musi to być wielokrotność NWW(12, 5) = 60. Najmniejsza trzycyfrowa: 120.}\par
3. 9 {\color{szary}– Bok musi dzielić 36 i 45: NWD = 9.}\par
4. A. 16 {\color{szary}– 30 : 16 daje resztę 14. Dzielniki 30: 1, 2, 3, 5, 6, 10, 15, 30.}\par
5. 12 {\color{szary}– Liczba zestawów musi dzielić 48 i 60: NWD(48, 60) = 12.}\par
6. 24 {\color{szary}– 1 + 2 + 7 + 14 = 24.}\par
7. 5 {\color{szary}– Wspólne dzielniki 15 i 35: 1, 5. Największy to 5.}\par
8. A. Prawda {\color{szary}– 24 = 2 · 12, więc co dzieli 12, dzieli też 24.}\par
9. 6 {\color{szary}– Pierwsza: 150, ostatnia: 900. To wielokrotności 150 od 1. do 6. – razem 6.}\par
10. B. Nie {\color{szary}– 1 + 2 + 3 + 4 + 6 = 16 $\neq$ 12.}\par
\medskip
\textbf{\color{pomd}Sprawdzian 3: Cechy podzielności}\par
1. B. Nie {\color{szary}– 1850 nie dzieli się przez 4 (dwie ostatnie cyfry: 50).}\par
2. 9450 {\color{szary}– Dzieli się przez 10, więc kończy się zerem; suma cyfr ma dzielić się przez 9: 9450 (suma 18).}\par
3. A. 4200 {\color{szary}– Liczba dzieli się przez 15, gdy dzieli się przez 3 i przez 5. 4200 = 15 · 280.}\par
4. 1 {\color{szary}– Reszta = 1, bo ostatnia cyfra 6.}\par
5. B. Fałsz {\color{szary}– Np. 6 jest parzysta.}\par
6. 5 {\color{szary}– Suma pozostałych cyfr to 13. Dopełniamy do wielokrotności 9: brakuje 5.}\par
7. A. Tak {\color{szary}– Wystarczy sprawdzić sumę cyfr: 2160 dzieli się przez 9 (suma jej cyfr: 9).}\par
8. 1 {\color{szary}– Suma pozostałych cyfr: 8. Najmniejsza cyfra, która dopełnia do wielokrotności 3, to 1.}\par
9. A. 25\,500 {\color{szary}– Cecha podzielności przez 10: ostatnia cyfra 0. 25\,500 ją spełnia.}\par
10. C. 2526 {\color{szary}– Liczba dzieli się przez 6, gdy dzieli się przez 2 i przez 3. 2526 = 6 · 421.}\par
\medskip
\textbf{\color{pomd}Sprawdzian 4: Liczby pierwsze}\par
1. C. 6 {\color{szary}– Liczba kończąca się na 6 dzieli się przez 2, więc jest złożona.}\par
2. A. Prawda {\color{szary}– Inne parzyste dzielą się przez 2.}\par
3. np. 13 {\color{szary}– Np. 13: 221 = 13 · 17.}\par
4. 83 {\color{szary}– Sprawdzamy kolejne liczby po 82: pierwszą liczbą pierwszą jest 83.}\par
5. 1 {\color{szary}– 59 ma dzielniki tylko 1 i 59, a 61 nie dzieli się przez 59, więc NWD = 1.}\par
6. C. 83 {\color{szary}– 83 ma tylko dwa dzielniki: 1 i 83. Pozostałe są złożone.}\par
7. 7 {\color{szary}– To: 31, 37, 41, 43, 47, 53, 59.}\par
8. D. 121 {\color{szary}– 121 = 11 · 11, więc ma więcej niż dwa dzielniki.}\par
9. 3 {\color{szary}– Dzielniki: 1, 7, 49 – razem 3.}\par
10. 5 {\color{szary}– To: 41, 43, 47, 53, 59.}\par
\medskip
\textbf{\color{pomd}Sprawdzian 5: Rozkład na czynniki}\par
1. 20 {\color{szary}– 2 · 2 · 5 = 20.}\par
2. 3 {\color{szary}– 384 = 2 · 2 · 2 · 2 · 2 · 2 · 2 · 3.}\par
3. 3 {\color{szary}– 45 = 3 · 3 · 5.}\par
4. 2 {\color{szary}– 298 = 2 · 149 – 2 czynniki.}\par
5. 180 {\color{szary}– NWW = 2 · 2 · 3 · 3 · 5 = 180.}\par
6. 4 {\color{szary}– Wspólne czynniki: 2 · 2 = 4.}\par
7. 90 {\color{szary}– NWW = 2 · 3 · 3 · 5 = 90.}\par
8. 6 {\color{szary}– Wspólne czynniki: 2 · 3 = 6.}\par
\medskip
\textbf{\color{pomd}Sprawdzian zbiorczy}\par
1. 1 {\color{szary}– 23 ma dzielniki tylko 1 i 23, a 191 nie dzieli się przez 23, więc NWD = 1.}\par
2. 0 {\color{szary}– Reszta = 0, bo suma cyfr 27 daje z dzielenia przez 9 resztę 0.}\par
3. A. Tak {\color{szary}– Wystarczy sprawdzić sumę cyfr: 9186 dzieli się przez 3 (suma jej cyfr: 24).}\par
4. 40 {\color{szary}– 40 to najmniejsza liczba (różna od 0), która dzieli się przez 5, 8, 10.}\par
5. 6 {\color{szary}– Liczba zestawów musi dzielić 18 i 24: NWD(18, 24) = 6.}\par
6. 560 {\color{szary}– 560 = 80 · 7; następna wielokrotność 640 jest już $\geq$ 632.}\par
7. D. 22 {\color{szary}– 40 : 22 daje resztę 18. Dzielniki 40: 1, 2, 4, 5, 8, 10, 20, 40.}\par
8. D. 42\,446 {\color{szary}– Cecha podzielności przez 2: ostatnia cyfra 0, 2, 4, 6 lub 8. 42\,446 ją spełnia.}\par
9. A. Prawda {\color{szary}– 1 i 13.}\par
10. 2 {\color{szary}– 56 = 2 · 2 · 2 · 7.}\par
11. 13 {\color{szary}– 5 · x − 3 · x = 2 · x = 26, więc x = 26 : 2 = 13.}\par
12. 41 {\color{szary}– Sprawdzamy kolejne liczby po 39: pierwszą liczbą pierwszą jest 41.}\par
13. B. Nie {\color{szary}– 1 + 2 + 3 + 6 + 9 = 21 $\neq$ 18.}\par
14. 4 {\color{szary}– 228 = 2 · 2 · 3 · 19 – 4 czynniki.}\par
15. 400 {\color{szary}– 400 = 40 · 10, a poprzednia wielokrotność 360 nie jest większa od 381.}\par
16. C. 169 {\color{szary}– 169 = 13 · 13, więc ma więcej niż dwa dzielniki.}\par
\medskip